% AIPL の型健全性と型安全性（第2版）
%   lualatex -interaction=nonstopmode aipl_soundness2_ja.tex   （2回）
\documentclass[11pt,a4paper]{ltjsarticle}
\usepackage{amsmath,amssymb}
\usepackage{amsthm}
\usepackage{mathpartir}
\usepackage{listings}
\usepackage{xcolor}
\usepackage{geometry}
\usepackage{array}
\usepackage{booktabs}
\usepackage{longtable}
\usepackage[hidelinks]{hyperref}
\geometry{margin=22mm}

\newtheorem{theorem}{定理}
\newtheorem{lemma}[theorem]{補題}
\newtheorem{corollary}[theorem]{系}
\theoremstyle{definition}
\newtheorem{definition}[theorem]{定義}
\theoremstyle{remark}
\newtheorem{remark}[theorem]{注意}
\newtheorem{finding}[theorem]{所見}

\definecolor{codebg}{gray}{0.96}
\lstset{
  basicstyle=\ttfamily\small, backgroundcolor=\color{codebg},
  frame=single, framesep=4pt, rulecolor=\color{gray},
  breaklines=true, showstringspaces=false,
  columns=fullflexible, keepspaces=true,
  keywordstyle=\bfseries, commentstyle=\itshape\color{gray!130},
}
\lstdefinelanguage{AIPL}{
  morekeywords={class,method,var,new,now,future,await,reply,send,
    timeout,else,if,while,do,select,case,self,sender,remote,print,
    int,float,string,bool,unit,any,true,false},
  sensitive=true, morecomment=[l]{//}, morestring=[b]",
}

% --- 記法（第1版を踏襲） ---
\newcommand{\tint}{\ensuremath{\mathtt{int}}}
\newcommand{\tbool}{\ensuremath{\mathtt{bool}}}
\newcommand{\tunit}{\ensuremath{\mathtt{unit}}}
\newcommand{\tstr}{\ensuremath{\mathtt{string}}}
\newcommand{\tact}[1]{\ensuremath{\mathtt{actor}[#1]}}
\newcommand{\tfut}[2]{\ensuremath{\mathtt{future}\langle #1 \,!\, #2 \rangle}}
\newcommand{\Om}{\Omega}
\newcommand{\Sm}{\Sigma}
\newcommand{\Ph}{\Phi}
\newcommand{\Ps}{\Psi}
\newcommand{\mtab}{\mathrm{mtab}}
\newcommand{\stype}{\mathrm{stype}}
\newcommand{\mret}{\mathrm{ret}}
\newcommand{\meff}{\mathrm{eff}}
\newcommand{\heapok}{\mathrm{heap\_ok}}
\newcommand{\msgok}{\mathrm{msg\_ok}}
\newcommand{\taskok}{\mathrm{task\_ok}}
\newcommand{\confok}{\mathrm{conf\_ok}}
\newcommand{\futok}{\mathrm{fut\_ok}}
\newcommand{\stuck}{\mathrm{stuck}}
\newcommand{\cstep}{\Longrightarrow}
\newcommand{\subs}[2]{[#1 := #2]}
\newcommand{\Eff}{\mathcal{E}}

\newcommand{\tfutl}[2]{\ensuremath{\mathtt{future}\langle #1 \rangle @ #2}}
\newcommand{\lvl}{\mathrm{lvl}}
\newcommand{\prodok}{\mathrm{prod\_ok}}

\title{AIPL の型健全性・型安全性・デッドロック自由（第 3 版）\\
{\large ---\ 待ちを取り除かずにデッドロック自由を証明する。
AIPL$^{-2}$ の機械検証 ---}}
\author{児玉工房 / AICE $\cdot$ AIPL プロジェクト}
\date{2026年08月23日}

\begin{document}
\maketitle

\begin{abstract}
第 1 版（2026 年 7 月 28 日）は、\texttt{reply} と戻り値型を結んだ部分言語
AIPL$^-$ について型健全性と型安全性を Coq で機械検証した。
ただし\textbf{デッドロック自由は \texttt{await} を言語から取り除いた断片
についてしか言えていなかった}。待てる構文を捨てて「待ちが返らない」を
消していたのだから、当然である。

第 2 版（同年 7 月 30 日）は、期限を必須にすることで進行定理を二択にした。
ただしこれは紙の上の証明であり、しかも得られる性質は
\emph{デッドロックが有限時間の失敗に変わる}ことであって、
\emph{詰まらない}ことではない。互いに待ち合う二者は両方が期限切れになる。

本稿（第 3 版）は\textbf{待ちを取り除かず、期限にも頼らず}、
デッドロック自由そのものを機械検証する。
道具は\textbf{義務レベル}である --- \texttt{future} の型に
「それを埋める側のレベル」を持たせ、待ちは必ずレベルの高い方へ向かうことを
型で要求する。加えて実行時の不変条件を一つ置く ---
\textbf{未解決の \texttt{future} には必ずそれを埋める者がいる}。
この二つから、全タスクが待ち状態になることはない。

主結果は\textbf{進行定理から枝が一つ消える}ことである。

\begin{center}
\begin{tabular}{ll}
\toprule
& 進行定理の主張 \\
\midrule
第 1 版 & 終状態 $\lor$ 一歩進める $\lor$ \textbf{全タスクが待ち} \\
第 3 版 & 終状態 $\lor$ 一歩進める \\
\bottomrule
\end{tabular}
\end{center}

併せて\textbf{効果}も形式化に入れた。\texttt{future} の型に
レベルと同じやり方で効果を載せ、\textbf{送信は引き継がず待ちが引き継ぐ}という
実装どおりの規則を与えると、効果の健全性が保存定理から出る ----
\textbf{メソッドが宣言した効果は、\texttt{now} で呼んだ先も含めて、
そのメソッドが実際に行うことを覆っている}。

定理は十三あり、いずれも \texttt{Print Assumptions} が
\texttt{Closed under the global context}（公理なし）である。
併せて、三つの性質（健全性・安全性・デッドロック自由）が同時に成り立つ
\textbf{最大の構文}を機構ごとに調べ、
\texttt{select}・\texttt{timeout}・\texttt{any}／多相・\texttt{become} を
入れると何が壊れるかを明記する。

定理が空虚でないことも確かめた。\texttt{await} と \texttt{while} を実際に使う
プログラムを形式化し、仮定をすべて満たすことを示し、
併せて\textbf{レベルを揃えると本体の型検査が通らない}ことを確かめた ---
条件が実際に効いていることの対照である。
\end{abstract}

\tableofcontents

% =====================================================================
\part{位置づけ}
\label{part:place}

\section{三つの版の関係}
\label{sec:versions}

\begin{center}
\small
\begin{tabular}{@{}lllp{5.2cm}@{}}
\toprule
版 & 日付 & 証明の形 & デッドロック自由の言い方 \\
\midrule
第 1 版 & 2026-07-28 & 機械検証（Coq） &
  \texttt{await} を\textbf{言語から除いた}断片についてのみ \\
第 2 版 & 2026-07-30 & 紙 &
  期限を必須にして、\textbf{有限時間の失敗}に変換する \\
\textbf{第 3 版（本稿）} & 2026-08-23 & \textbf{機械検証（Coq）} &
  \texttt{await} を含んだまま、\textbf{そもそも詰まらない}ことを言う。
  効果も機械検証に入れた \\
\bottomrule
\end{tabular}
\end{center}

三つは重なっているが同じではない。違いは\textbf{何によって詰まらないと言うか}である。

\begin{itemize}
\item 第 1 版は\textbf{待つ構文を捨てた}。待たなければ詰まらない。
\item 第 2 版は\textbf{待ちに上限を置いた}。上限があれば有限時間で終わる。
      ただし終わり方が失敗でもよい。
\item 第 3 版は\textbf{待ちの向きを揃えた}。循環しなければ詰まらない。
      失敗に変換する必要が無い。
\end{itemize}

\begin{remark}[名前について]
第 2 版も、そこで切り出した部分言語を AIPL$^{-2}$ と呼んでいる。
本稿の AIPL$^{-2}$ とは中身が違う --- 第 2 版のそれは
\textbf{効果と期限}を持ち、本稿のそれは\textbf{義務レベル}を持つ。
名前が衝突しているのは、どちらも「第 1 版の次」を意味していたためである。

本稿では、\textbf{機械検証された言語}を AIPL$^{-2}$ と呼ぶ。
第 2 版の言語を区別して指す必要があるときは AIPL$^{-\mathcal{E}}$ と書く。

なお本稿の執筆中に\textbf{効果を機械検証へ合流させた}（\S\ref{sec:effects}）。
残る隔たりは\textbf{期限}だけである。
\end{remark}

\section{何を証明したいのか}

三つある。

\begin{center}
\small
\begin{tabular}{@{}lp{9.4cm}@{}}
\toprule
性質 & 主張 \\
\midrule
\textbf{型健全性} & 型の付いた構成が一歩進んでも、型の付いた構成のままである
  （保存）。加えて、各アクターの状態と解決済み \texttt{future} の値は
  常に宣言された型を持つ \\
\textbf{型安全性} & 型の付いた構成から到達できる構成は $\stuck$ でない。
  「終わってもいないのに一歩も進めない」ことが起きない \\
\textbf{デッドロック自由} & 型の付いた構成が \textbf{blocked} にならない。
  すなわち「メールボックスが空で、全タスクが待ち状態」という構成が現れない \\
\bottomrule
\end{tabular}
\end{center}

三つ目が本稿の主題である。第 1 版はこれを、待てる構文を除いた断片について
しか言えなかった。

% =====================================================================
\part{定式化}
\label{part:form}

\section{構文}
\label{sec:syntax}

型と項は次のとおりである。実装の \texttt{src/ast.ml} と
\texttt{src/types.ml} のアクター核を取り出したものである。

\[
\tau ::= \tint \mid \tbool \mid \tunit \mid \tact{c} \mid \tfutl{\tau}{n}
\]

\[
\begin{array}{rcl}
e &::=& n \mid b \mid () \mid x \mid \mathtt{self}
      \mid \mathtt{o}_o \mid \mathtt{f}_k \\
  &\mid& e + e \mid e < e \mid \mathtt{if}\ e\ e\ e
      \mid \mathtt{var}\ x = e;\ e \\
  &\mid& \mathtt{get} \mid \mathtt{set}\ e \mid \mathtt{new}\ c \\
  &\mid& \mathtt{future}\ e.m(e) \mid \mathtt{await}\ e \\
  &\mid& e ; e \mid \mathtt{while}\ e\ \mathtt{do}\ e
\end{array}
\]

$\mathtt{o}_o$ と $\mathtt{f}_k$ は実行時にだけ現れる参照である。
$\tfutl{\tau}{n}$ が本稿の要で、
「型 $\tau$ の値を、\textbf{レベル $n$ のメソッドが}埋める \texttt{future}」を表す。

\subsection{糖衣}

次の三つは規則を増やさずに書ける。

\begin{center}
\begin{tabular}{@{}ll@{}}
\toprule
書き方 & 展開 \\
\midrule
\texttt{now t.m(e)} & $\mathtt{await}\ (\mathtt{future}\ t.m(e))$ \\
\texttt{send t.m(e);} & $(\mathtt{future}\ t.m(e))\, ;\, ()$ \\
\texttt{e1; e2} & 逐次実行（構文として持つ） \\
\bottomrule
\end{tabular}
\end{center}

\texttt{while} は \texttt{if} へ展開する形にした。
束縛を作らないので変数捕獲の心配が無い。

\[
\mathtt{while}\ c\ \mathtt{do}\ b
\;\longrightarrow\;
\mathtt{if}\ c\ (b\, ;\, \mathtt{while}\ c\ \mathtt{do}\ b)\ ()
\]

\section{表と構成}

プログラムは五つの表で与える。

\begin{center}
\small
\begin{tabular}{@{}ll@{}}
\toprule
表 & 内容 \\
\midrule
$\stype(c)$ & クラス $c$ の状態の型 \\
$\mathrm{sinit}(c)$ & クラス $c$ の状態の初期値 \\
$\mtab(c,m)$ & $c.m$ の引数型と返り値型 \\
$\mathrm{body}(c,m)$ & $c.m$ の本体（引数は変数 $0$） \\
$\lvl(c,m)$ & \textbf{$c.m$ の義務レベル} \\
\bottomrule
\end{tabular}
\end{center}

実行時の構成は $(\Om, \mathit{ms}, \mathit{ts})$ である。
$\Om$ はヒープで、アクター表・状態表・\texttt{future} 表・\texttt{future} の値を持つ。
\texttt{future} 表は $(\tau, n)$ の対を持つ --- 型と、\textbf{埋める側のレベル}である。
$\mathit{ms}$ はメールボックス（飛んでいるメッセージの並び）、
$\mathit{ts}$ は走っているタスクの並びである。
タスクは $(o, k, e)$ の三つ組で、
「アクター $o$ の上で式 $e$ を評価し、その値で \texttt{future} $k$ を埋める」を意味する。

\section{型付け}
\label{sec:typing}

判断は
\[
\Om;\Ph \vdash_{c,L}\ \Gamma \vdash e : \tau
\]
と読む。$\Om$ はアクター表、$\Ph$ は \texttt{future} 表、
$c$ は\textbf{いま検査しているメソッドが属するクラス}、
$L$ は\textbf{そのメソッドの義務レベル}である。
$L$ を判断に持たせたのが第 1 版との違いである。

主要な規則だけ挙げる。

\begin{center}
\small
\begin{tabular}{@{}lp{9.0cm}@{}}
\toprule
規則 & 内容 \\
\midrule
参照 & $\Ph(k) = (\tau, n)$ ならば $\mathtt{f}_k : \tfutl{\tau}{n}$ \\
\addlinespace
送信 & $e_0 : \tact{c'}$ かつ $\mtab(c',m) = (\tau_a,\tau_r)$ かつ $e_1 : \tau_a$
  ならば\newline
  $\mathtt{future}\ e_0.m(e_1) : \tfutl{\tau_r}{\lvl(c',m)}$ \\
\addlinespace
\textbf{待ち} & $e : \tfutl{\tau}{n}$ かつ $\boldsymbol{L < n}$ ならば
  $\mathtt{await}\ e : \tau$ \\
\bottomrule
\end{tabular}
\end{center}

\begin{center}
\fbox{\parbox{0.93\textwidth}{
\textbf{待ちは必ず上へ向かう。} レベル $L$ のメソッドが待てるのは、
レベルが $L$ より\emph{真に大きい}メソッドが埋める \texttt{future} だけである。
これがデッドロック自由の全部である。}}
\end{center}

プログラム全体が型検査を通っているとは、次を満たすことである。

\begin{center}
\small
\begin{tabular}{@{}lp{8.4cm}@{}}
\toprule
仮定 & 内容 \\
\midrule
$\mathrm{sinit\_value}$ & 状態の初期値は値である \\
$\mathrm{sinit\_ok}$ & 状態の初期値は宣言された型を持つ \\
$\mathrm{bodies\_ok}$ & $\mtab(c,m) = (\tau_a,\tau_r)$ ならば、
  $c.m$ の本体は\textbf{レベル $\lvl(c,m)$ のもとで} $\tau_r$ の型を持つ \\
\bottomrule
\end{tabular}
\end{center}

これらは Coq の \texttt{Section} の \texttt{Hypothesis} であり、
\texttt{End} で各定理の前提に変わる。\textbf{公理ではない}。

\section{操作的意味論}

タスク一歩の関係 $\Om; o \vdash e \to \Om'; \mathit{out}; e'$ と、
構成一歩の関係 $\cstep$ の二段である。$\mathit{out}$ は新たに送出される
メッセージの並びである。

構成の一歩は三つしかない。

\begin{center}
\small
\begin{tabular}{@{}lp{8.8cm}@{}}
\toprule
規則 & 内容 \\
\midrule
CTask & タスクが一歩進む。送出されたメッセージはメールボックスに入る \\
CFinish & タスクの式が値になった。その値が \texttt{reply} の値であり、
  担当していた \texttt{future} を埋めてタスクは消える \\
CDeliver & メッセージを配送し、メソッド本体を新しいタスクとして起こす \\
\bottomrule
\end{tabular}
\end{center}

送信の規則が要点である。

\[
\Om; o \vdash \mathtt{future}\ \mathtt{o}_{o'}.m(v)
\;\to\;
\Om[\Ph \mathbin{+\!\!+} (\tau_r, \lvl(c',m))];\;
[(o', m, v, |\Ph|)];\;
\mathtt{f}_{|\Ph|}
\]

すなわち\textbf{\texttt{future} を一つ増やすのと同時に、
それを埋めるためのメッセージを必ず一通出す}。
この対応が次節の不変条件になる。

\section{不変条件}
\label{sec:inv}

構成が正しいとは、次の五つを満たすことである。

\begin{center}
\small
\begin{tabular}{@{}lp{8.6cm}@{}}
\toprule
条件 & 内容 \\
\midrule
$\heapok$ & 表の長さが揃い、各アクターの状態は宣言された型を持ち、
  解決済み \texttt{future} の値は記録された型を持つ \\
$\msgok$ & 飛んでいるメッセージの宛先は実在し、そのクラスはそのメソッドを持ち、
  引数の型が合い、\textbf{担当する \texttt{future} のレベルが $\lvl(c,m)$ に一致する} \\
$\taskok$ & タスクは、\textbf{自分が埋める \texttt{future} のレベルのもとで}型が付く \\
$\prodok$ & \textbf{未解決の \texttt{future} には必ず「埋める者」がいる} ---
  メールボックスの中のメッセージか、走っているタスクのどちらかである \\
$\mathrm{ext}$ & 起動時のアクター表は現在のアクター表に含まれる \\
\bottomrule
\end{tabular}
\end{center}

$\prodok$ が第 3 版で足したものである。
送信が \texttt{future} とメッセージを同時に作るので初期状態で成り立ち、
配送はメッセージをタスクに変え（担い手が移る）、
完了は \texttt{future} を埋める（義務が消える）ので、一歩ごとに保たれる。

% =====================================================================
\part{証明}
\label{part:proof}

\section{主定理}

\begin{center}
\small
\begin{tabular}{@{}lp{8.4cm}@{}}
\toprule
定理 & 主張 \\
\midrule
\texttt{no\_method\_not\_understood} &
  型の付いた構成を飛ぶメッセージは、宛先が実在し、そのクラスがそのメソッドを持ち、
  引数と \texttt{future} の型もレベルも合っている \\
\texttt{preservation} & 一歩進んでも $\confok$ が保たれる \\
\texttt{preservation\_star} & 何歩進んでも保たれる \\
\texttt{type\_safety} & 到達できる構成は $\stuck$ でない \\
\textbf{\texttt{deadlock\_free}} & 型の付いた構成は blocked にならない \\
\textbf{\texttt{progress\_total}} & 終状態か、一歩進めるか、の二択 \\
\texttt{deadlock\_free\_star} & 到達できる構成すべてについて同じことを言う \\
\texttt{state\_type\_invariant} & どの時点でも各アクターの状態は宣言型を持つ \\
\texttt{future\_type\_invariant} & 解決済み \texttt{future} の値は宣言された返り値型を持つ \\
\addlinespace
\texttt{effect\_no\_increase} & 一歩進んでも効果は増えない \\
\texttt{effect\_soundness} & タスクの効果は宣言された効果に収まり続ける \\
\texttt{await\_charges\_callee} & 待つと呼び先の効果を必ず引き継ぐ \\
\texttt{send\_does\_not\_charge} & 送るだけでは引き継がない \\
\bottomrule
\end{tabular}
\end{center}

\section{デッドロック自由の証明}
\label{sec:df}

\begin{theorem}[デッドロック自由]
$\confok(C)$ ならば $C$ は blocked でない。
\end{theorem}

\begin{proof}
$C = (\Om, \mathit{ms}, \mathit{ts})$ が blocked だと仮定する。
定義より $\mathit{ms} = []$、$\mathit{ts} \neq []$、
そして $\mathit{ts}$ のすべてのタスクが待ち状態である。

タスクの\textbf{レベル}を、そのタスクが埋める \texttt{future} に
記録されているレベルと定める。
$\mathit{ts}$ は空でない有限の並びなので、
\textbf{レベルが最大のタスク} $t = (o,k,e)$ が取れる。

$t$ も待ち状態である。待っている式は必ず未解決の \texttt{future} を
ひとつ名指ししており、$\taskok$ よりタスクはレベル $L = \lvl(k)$ のもとで
型が付いているから、待ちの規則によりその \texttt{future} のレベル $n$ は
$L < n$ を満たす。

$\prodok$ より、その \texttt{future} には埋める者がいる。
$\mathit{ms} = []$ だからメッセージではありえず、したがってタスクである。
そのタスクが埋める \texttt{future} のレベルは $n$ であり、$n > L$ である。
これは $t$ のレベル最大性に反する。
\end{proof}

\begin{corollary}[進行、二択]
$\confok(C)$ ならば、$C$ は終状態であるか、一歩進める。
\end{corollary}

\begin{proof}
第 1 版の進行定理は三択（終状態・一歩進める・blocked）である。
上の定理が第三の枝を消す。
\end{proof}

\begin{finding}
証明に効いているのは二つだけである --- \textbf{待ちの規則}（$L < n$）と、
\textbf{$\prodok$}（未解決の \texttt{future} には埋める者がいる）である。
前者は型の側、後者は実行時の側にある。
片方だけでは足りない。
待ちの向きが揃っていても、埋める者が居なければ待ちは返らない。
埋める者が居ても、向きが循環していれば互いを待つ。
\textbf{型と実行時の不変条件が噛み合って初めて言える}性質である。
\end{finding}

\section{公理を使っていないこと}

九つの定理すべてについて \texttt{Print Assumptions} を実行し、
\texttt{Closed under the global context} を確認した。
\texttt{Makefile} の \texttt{check} ターゲットがこれを自動で行う。

% =====================================================================
\part{範囲}
\label{part:scope}

\section{証明できる最大の構文}
\label{sec:max}

三つの性質が\emph{同時に}成り立つ範囲はどこまでか。
機構ごとに調べた。

\subsection{入れても保たれるもの}

\begin{center}
\small
\begin{tabular}{@{}lp{8.6cm}@{}}
\toprule
機構 & 理由 \\
\midrule
\texttt{while} / 逐次実行 &
  止まらない繰り返しは「デッドロック」ではなく「停止しない」である。
  進行定理は一歩進めることしか言っていないので、三つとも保たれる。
  \texttt{while} は \texttt{if} へ展開する形にしたので束縛を作らず、
  変数捕獲が起きない \\
動的なアクター生成 & \texttt{new} はヒープを伸ばすだけである \\
\texttt{future} の入れ子 &
  \texttt{await} の中に \texttt{await} があってもよい。
  レベルは式の位置ではなく\textbf{メソッド}に付くので、
  入れ子でも同じ条件で足りる \\
\bottomrule
\end{tabular}
\end{center}

\subsection{入れると壊れるもの}

\begin{center}
\small
\begin{tabular}{@{}lp{8.6cm}@{}}
\toprule
機構 & 壊れるもの \\
\midrule
\texttt{select} &
  待つ相手が誰かは型に現れないので、レベルが効かない。
  進行の枝が一つ増え、デッドロック自由が壊れる。
  実装では期限の義務づけと送り手の存在検査で補っている \\
\texttt{timeout} &
  時間を意味論に入れる必要がある。進行は保てるが、
  主張が「詰まらない」から「詰まっても諦める」に変わる ---
  第 2 版がまさにこれである \\
\texttt{any} / 多重定義 / 多相 &
  標準形補題が壊れる。値の形が型から決まらないので進行が通らない \\
\texttt{become} &
  受け手のクラスが実行時に変わると、配送時に型が合っている保証が崩れる
  （\texttt{no\_method\_not\_understood} が壊れる）。
  なお実装では \texttt{become} 自体を言語から外した \\
他アクターの状態への直接代入 &
  $\heapok$ の「各アクターの状態は宣言型を持つ」が、
  自分のメソッド内でしか保てなくなる \\
\bottomrule
\end{tabular}
\end{center}

\begin{finding}
この線引きは、言語の設計として機構を採るか落とすかを決めた判断とは
\textbf{独立に}引いたのに、ほぼ同じところに落ちた。
\texttt{become} も \texttt{any} も多相も、設計として落とした理由と、
証明が通らない理由が一致している。
\textbf{証明の通る範囲は、設計判断の後追いではなく、その裏づけであった}。
\end{finding}

\section{空虚でないことの確認}
\label{sec:nonvacuous}

定理はすべて「プログラムが型検査を通っている」という仮定のもとにある。
レベルの条件が満たせないなら定理は空虚なので、確かめた。

\begin{lstlisting}[language=AIPL]
class Svc    { method get(x) : int { reply(x); } }                  // レベル 1
class Caller { method run(x) : int {                                // レベル 0
    while (x < 0) { }                    // 第3版で広げた構文
    reply(now svc.get(x));               // await を含む
} }
\end{lstlisting}

このプログラムについて、$\mathrm{sinit\_value}$・$\mathrm{sinit\_ok}$・
$\mathrm{bodies\_ok}$ をすべて Coq で証明し、
初期構成が $\confok$ であることを示した。
したがってデッドロック自由が適用でき、進行も二択で言える。

\textbf{対照実験}も行った。$\lvl(\mathrm{Svc},\mathrm{get})$ を $1$ から $0$ に
変えると（すなわち待ちの向きを平らにすると）、
$\mathrm{bodies\_ok}$ の証明が通らなくなる ---
$0 < 0$ が示せないからである。条件が実際に効いている。

\section{効果}
\label{sec:effects}

レベルを載せたのと\textbf{同じやり方}で効果も載せた。

\subsection{型に載せる}

\[
\tfutl{\tau}{n} \;\longrightarrow\;
\mathtt{future}\langle \tau \rangle @ n \,!\, \varepsilon
\]

判断も一つ増える。

\[
\Om;\Ph \vdash_{c,L}\ \Gamma \vdash e : \tau \,!\, \varepsilon
\]

$\varepsilon$ は効果名の集合である（形式化では \texttt{list nat}、
包含は \texttt{incl}、合併は \texttt{++}）。

\subsection{規則の要は二つ}

\begin{center}
\small
\begin{tabular}{@{}lp{8.6cm}@{}}
\toprule
規則 & 効果 \\
\midrule
送信 & $\mathtt{future}\ e_0.m(e_1)$ の効果は
  $\varepsilon_0 \cup \varepsilon_1$ である ----
  \textbf{呼び先の効果を引き継がない}。待たないからである。
  ただし結果の型には $\meff(c',m)$ を\emph{記録する} \\
\addlinespace
\textbf{待ち} & $\mathtt{await}\ e$ の効果は
  $\varepsilon_e \cup \varepsilon_c$ である ----
  \textbf{呼び先の効果を引き継ぐ}。$\varepsilon_c$ は
  \texttt{future} の型に記録されていたものである \\
\bottomrule
\end{tabular}
\end{center}

これは実装の挙動と同じである。実装でも \texttt{send} は伝播せず、
\texttt{await} が伝播する。
プログラム全体の仮定にも一つ足す ----
$\mathrm{bodies\_ok}$ は「本体の効果が宣言した効果 $\meff(c,m)$ に収まる」を
要求する。実装が推論した効果を注釈と照合するのと同じ形である。

\subsection{不変条件}

$\taskok$ に一項足す ---- タスクの効果は、
そのタスクが埋める \texttt{future} に記録された効果に収まる。
配送のときに $\mathrm{bodies\_ok}$ が入り、以後は保存で保たれる。

\subsection{定理}

\begin{center}
\small
\begin{tabular}{@{}lp{8.0cm}@{}}
\toprule
定理 & 主張 \\
\midrule
\texttt{effect\_no\_increase} & 一歩進んでも効果は増えない \\
\texttt{effect\_soundness} & 走っているタスクの効果は、担当する
  \texttt{future} に記録された効果に収まり続ける \\
\texttt{await\_charges\_callee} & 待つと呼び先の効果を必ず引き継ぐ \\
\texttt{send\_does\_not\_charge} & 送るだけでは引き継がない \\
\bottomrule
\end{tabular}
\end{center}

\texttt{effect\_soundness} の中身は次のとおりである ----
\textbf{メソッドが宣言した効果は、そのメソッドが実際に行うことを
（\texttt{now} で呼んだ先も含めて）覆っている}。

\begin{corollary}[一段隔てても隠せない]
$\meff(c,m)$ に $\mathtt{ai}$ が入っていなければ、
$c.m$ の本体は $\mathtt{ai}$ を持つメソッドを \texttt{await} できない。
\end{corollary}

\begin{proof}
待ちの規則より、$\mathtt{await}$ の効果は呼び先の効果を含む。
$\mathrm{bodies\_ok}$ より本体の効果は $\meff(c,m)$ に収まる。
包含の推移性による。
\end{proof}

これが実装の動機そのものである ----
\textbf{実時間で動くアクターが、機外のモデルを同期で待てない}。

\begin{finding}
効果を足すのに、レベルのときと同じ形が使えた ----
\textbf{\texttt{future} の型に載せ、待ちで課金する}。
違いは「レベルは比較し、効果は合併する」ことだけである。
どちらも「呼び先について知りたいことを、\texttt{future} の型に持たせる」
という一つの考え方の例である。
\end{finding}

\section{言えないこと}
\label{sec:cannot}

\begin{itemize}
\item \textbf{停止性}。デッドロック自由は「詰まらない」であって
  「必ず終わる」ではない。\texttt{while}\ $1>0$\ \texttt{do}\ $\{\}$ は
  止まらないが、それは待ちではないので進行定理と矛盾しない。
\item \textbf{飢餓}。あるタスクが永久に選ばれない可能性は排除していない。
  進行定理は「どれかが進める」しか言わない。
\item \textbf{宛先が動的な場合}。レベルは $\mtab$ から引くので、
  宛先のクラスが静的に決まらなければ効かない。
\item \textbf{同じクラスの別インスタンス}。レベルはメソッド単位なので、
  親が子を待つ木構造のような場合も一律に扱われる。
  緩めるにはレベル変数が要る。
\item \textbf{期限}。本稿の形式化には入っていない（次節）。
  効果は入れた（\S\ref{sec:effects}）。
\end{itemize}

\section{残っていること}
\label{sec:remains}

\subsection{効果は入れた}

本稿の執筆中に足した（\S\ref{sec:effects}）。
第 2 版が紙で示していた効果の健全性が、機械検証になった。

\subsection{期限を足す}

期限を入れると時間を意味論に持ち込む必要がある。
得られるのは「詰まっても有限時間で失敗する」であって、
本稿の「詰まらない」より弱い。
両方を持つ形（レベルで循環を止め、期限で残りを救う）が実装の姿なので、
形式化としてもそこが終着点になる。

\subsection{名前の整理}

効果が入ったので、第 2 版の AIPL$^{-\mathcal{E}}$ と本稿の AIPL$^{-2}$ の
違いは\textbf{期限だけ}になった。期限を足した時点で名前は一本化できる。
それまでは本稿の記法（AIPL$^{-2}$ は機械検証された言語、
AIPL$^{-\mathcal{E}}$ は第 2 版の言語）で区別する。

% =====================================================================
\section{まとめ}

第 1 版はデッドロック自由を、\texttt{await} を言語から取り除いた断片に
ついてしか言えなかった。第 2 版は期限で有限時間の失敗に変換した。
本稿は\textbf{待ちを取り除かず、期限にも頼らず}、
デッドロック自由そのものを機械検証した。

足したのは四つだけである ---
\texttt{future} の型にレベルを持たせ、判断にレベルを足し、
待ちに $L < n$ を要求し、
「未解決の \texttt{future} には埋める者がいる」を不変条件に置いた。
証明の骨は「レベルが最大のタスクを取る」ことである。

そして三つの性質が同時に成り立つ最大の構文を調べたところ、
その境界は\textbf{言語の設計判断と独立に引いたのに、ほぼ同じところに落ちた}。
証明の通る範囲は、設計の後追いではなく、その裏づけであった。

\subsection*{再現方法}

\begin{lstlisting}[language=]
cd ~/aios/abclcp/coq
make            # AIPLSoundness2.v と AIPLSoundness2Example.v を含めてビルド
make check      # 主定理の Print Assumptions を確認
\end{lstlisting}

\begin{center}
\small
\begin{tabular}{@{}ll@{}}
\toprule
ファイル & 内容 \\
\midrule
\texttt{coq/AIPLSoundness2.v} & 本稿の形式化と証明（1{,}314 行） \\
\texttt{coq/AIPLSoundness2Example.v} & 空虚でないことの確認（123 行） \\
\texttt{coq/AIPLSoundness.v} & 第 1 版（1{,}121 行） \\
\texttt{docs/aipl\_soundness2\_ja.pdf} & 第 2 版（紙の証明。効果と期限） \\
\texttt{docs/aipl\_paper5\_ja.pdf} & 論文。言語全体の設計と採否 \\
\bottomrule
\end{tabular}
\end{center}

\end{document}
